\section*{附录}
\addcontentsline{toc}{section}{附录}

\keypoint{附录只保留复现路径、核心模型实现和关键检验代码。完整 Python/R 程序随项目文件提交，避免大段源码削弱论文主体可读性。}

\subsection*{A 代码复现说明}

本文的建模、校验、绘图和论文生成均保存在同一项目目录中。Python 环境采用 \texttt{python=3.11}，主要依赖包括 \texttt{pandas}、\texttt{numpy}、\texttt{scipy}、\texttt{matplotlib}、\texttt{statsmodels}、\texttt{scikit-learn}、\texttt{plotly} 和 \texttt{pyyaml}。核心运行顺序如下。

\begin{paperlisting}{lst:reproduce-flow}{python}{论文结果复现主流程}
# 1. 清洗赛题附件数据
python -m src.pipeline.clean_data

# 2. 建立短期动态递推模型
python -m src.models.dynamic_short_term

# 3. 校准综合机制参数
python -m src.calibration.calibrate_dynamic_model

# 4. 生成 60--180 天三情景预测
python -m src.scenarios.scenario_forecast

# 5. 运行蒙特卡洛情景树与统计审计
python -m src.analysis.monte_carlo_scenario_tree
python -m src.analysis.statistical_audit

# 6. 生成论文图表并构建最终论文
python -m src.visualization.final_paper_figures
./scripts/build_final_paper.sh
\end{paperlisting}

项目核心结构如下。未列出的缓存文件、临时渲染文件和历史实验文件不影响最终复现。

{\footnotesize
\begin{verbatim}
数模/
|-- data/                 原始数据、清洗数据和外部校验数据
|-- src/
|   |-- pipeline/         数据清洗与校验
|   |-- models/           传统供需基准与短期动态模型
|   |-- calibration/      短期模型参数搜索与评价
|   |-- scenarios/        60--180 天长期三情景预测
|   |-- analysis/         敏感性、稳健性、GPR、OVX 与统计检验
|   |-- visualization/    论文正式图表生成
|-- scripts/              一键构建、R 审计与 DOCX 后处理脚本
|-- paper/                总论文 LaTeX 源码、章节和正式图表
|-- output/               模型结果、校验报告和最终 PDF/DOCX
\end{verbatim}
}

\subsection*{B 数据口径}

本文真实价格数据只来自赛题附件，并经清洗后形成中文命名数据文件。短期拟合使用附件中 2026 年 3 月 3 日至 2026 年 5 月 5 日的冲突窗口；2017--2026 年完整价格序列仅用于历史同长度窗口、基准模型分布和事件极端性检验。EIA 周度 SPR、商业库存和原油产量数据、JODI 多国月度产量、库存、进出口、需求和炼厂数据、OPEC 全球供需平衡表、GPR 月度指数和 OVX 隐含波动率用于外部合理性校验，不反向替代附件真实价格。

\subsection*{C 综合机制递推模型核心代码}

代码~\ref{lst:dynamic-recursion} 展示短期综合机制递推模型的核心实现。模型逐日计算有效需求、有效供应、剩余缺口、风险溢价和缓冲折价，并由目标价格递推得到模拟价格。

\begin{paperlisting}{lst:dynamic-recursion}{python}{综合机制递推模型核心实现}
def simulate_dynamic_model(event_df, assumptions, behavior):
    base_price = float(event_df.iloc[0]["pre_close"])
    previous_price = base_price
    previous_fear_excess = assumptions.fear_initial
    inventory_remaining = assumptions.commercial_inventory
    gap_closure_day = None
    rows = []

    for step_index, actual_row in enumerate(event_df.itertuples(index=False)):
        day_index = int((actual_row.trade_date - event_df["trade_date"].min()).days)
        elasticity = interpolate_elasticity(day_index, assumptions)
        price_ratio = max(previous_price / base_price, 0.1)

        # 需求侧：价格弹性与高油价需求收缩共同作用
        price_adjusted_demand = assumptions.base_demand * (price_ratio ** elasticity)
        demand_decline = ramp(day_index, 0, assumptions.demand_decline_ramp_days,
                              assumptions.observed_demand_decline)
        effective_demand = max(price_adjusted_demand - demand_decline,
                               assumptions.base_demand * 0.70)

        # 供给侧：基础供应扣除中断后，加入 SPR、绕道运输和库存缓冲
        spr_release = ramp(day_index, assumptions.spr_delay_days,
                           assumptions.spr_ramp_days, assumptions.spr_max_release)
        route_supply = ramp(day_index, assumptions.route_start_day,
                            assumptions.route_ramp_days, assumptions.route_max_capacity)
        supply_without_inventory = (assumptions.base_supply
                                    - assumptions.supply_interruption
                                    + spr_release + route_supply)

        raw_gap = max(effective_demand - supply_without_inventory, 0.0)
        inventory_buffer = min(raw_gap * behavior.inventory_response,
                               assumptions.inventory_daily_cap,
                               inventory_remaining)
        inventory_remaining -= inventory_buffer
        effective_supply = supply_without_inventory + inventory_buffer
        residual_gap = max(effective_demand - effective_supply, 0.0)
        if gap_closure_day is None and residual_gap <= assumptions.base_demand * 0.005:
            gap_closure_day = day_index

        # 价格形成层：物理缺口、地缘风险、恐慌溢价与市场缓冲折价
        fear_excess = assumptions.fear_initial * np.exp(-assumptions.fear_decay * day_index)
        shortage_pressure = (base_price * behavior.pressure_scale
                             * (residual_gap / assumptions.base_demand)
                             / max(abs(elasticity), 0.01))
        blockade_risk_premium = (base_price * behavior.risk_weight
                                 * (assumptions.supply_interruption / assumptions.base_demand)
                                 * (1 - np.exp(-day_index / 7))
                                 * np.exp(-BLOCKADE_RISK_DECAY * day_index))
        panic_premium = base_price * 0.45 * fear_excess

        target_price = (base_price + shortage_pressure + blockade_risk_premium
                        + panic_premium
                        - buffer_confirmation_discount(day_index, gap_closure_day,
                                                       base_price, behavior)
                        - expectation_relief_discount(day_index, base_price, behavior))

        simulated_price = previous_price + behavior.adjustment_speed * (target_price - previous_price)
        simulated_price += 2.5 * (fear_excess - previous_fear_excess)
        rows.append({"day_index": day_index, "simulated_price": simulated_price})
        previous_price = simulated_price
        previous_fear_excess = fear_excess

    return pd.DataFrame(rows)
\end{paperlisting}

\subsection*{D 参数校准模块}

代码~\ref{lst:calibration} 给出参数校准流程。本文不是手工指定最优参数，而是在固定机制结构下，综合使用差分进化、局部稳定性搜索和拟合质量精修选择代表性参数。

\begin{paperlisting}{lst:calibration}{python}{参数校准与局部稳健性搜索}
def refine_with_continuous_search(event_df, base_assumptions):
    def objective(values):
        assumptions, behavior = decode_continuous_parameters(values, base_assumptions)
        simulation = dynamic.simulate_dynamic_model(event_df, assumptions, behavior)
        metrics = evaluate_simulation(simulation)

        # 综合得分同时考虑整体误差、峰值误差、末日误差和平台解释能力
        return metrics["综合得分"] + excellence_penalty(metrics)

    result = differential_evolution(
        objective,
        CONTINUOUS_PARAMETER_BOUNDS,
        seed=RANDOM_SEED,
        maxiter=LOCAL_REFINEMENT_MAXITER,
        popsize=LOCAL_REFINEMENT_POPSIZE,
        polish=True,
        workers=1,
        tol=0.002,
    )

    assumptions, behavior = decode_continuous_parameters(result.x, base_assumptions)
    simulation = dynamic.simulate_dynamic_model(event_df, assumptions, behavior)
    metrics = evaluate_simulation(simulation)
    metrics["局部精修目标值"] = float(result.fun)
    return assumptions, behavior, simulation, metrics

def calibrate(event_df, base_assumptions):
    rows, simulations = [], {}
    for candidate_id, (assumptions, behavior) in enumerate(sampled_parameter_sets(base_assumptions)):
        simulation = dynamic.simulate_dynamic_model(event_df, assumptions, behavior)
        metrics = evaluate_simulation(simulation)
        rows.append({"candidate_id": candidate_id, "candidate_source": "seeded_random",
                     **metrics})
        simulations[candidate_id] = simulation

    # 在随机搜索基础上继续做连续精修、局部稳健性搜索和拟合质量精修
    refined = refine_with_continuous_search(event_df, base_assumptions)
    stable = refine_with_local_stability_search(event_df, base_assumptions,
                                                refined[0], refined[1])
    final = refine_with_fit_quality_search(event_df, base_assumptions,
                                           stable[0], stable[1])
    return select_representative_candidates(rows, simulations, final)
\end{paperlisting}

\subsection*{E 长期情景递推与政策缓冲}

代码~\ref{lst:future-recursion} 展示长期外推中 SPR 自适应释放与供需再平衡的处理。该模块避免把战略储备释放机械设为满额，而是根据剩余缺口、价格压力和时间衰减动态收缩释放强度。

\begin{paperlisting}{lst:future-recursion}{python}{长期情景递推中的 SPR 自适应释放}
def adaptive_spr_release(day_index, scheduled_spr, gross_gap_before_spr,
                         previous_price, base_price, assumptions):
    if scheduled_spr <= 0:
        return 0.0, 0.0

    # 当物理缺口已基本覆盖时，SPR 不再机械满额释放
    reserve_buffer = assumptions.base_demand * GAP_CLOSURE_SHARE
    coverage_need = max(gross_gap_before_spr + reserve_buffer, 0.0)
    demand_based_release = min(scheduled_spr, coverage_need)

    stress_ratio = np.clip(gross_gap_before_spr / max(assumptions.supply_interruption, 1.0),
                           0.0, 1.0)
    price_stress = np.clip((previous_price / base_price - SPR_PRICE_STRESS_START_RATIO)
                           / SPR_PRICE_STRESS_WIDTH, 0.0, 1.0)
    minimum_policy_share = SPR_MIN_POLICY_SHARE + SPR_PRICE_POLICY_SHARE * price_stress
    stress_share = max(minimum_policy_share, stress_ratio)

    time_taper = 1.0
    if day_index > SPR_TAPER_START_DAY:
        time_taper = SPR_TAPER_FLOOR + (1 - SPR_TAPER_FLOOR) * np.exp(
            -(day_index - SPR_TAPER_START_DAY) / SPR_TAPER_DECAY_DAYS
        )

    actual_release = max(demand_based_release, scheduled_spr * stress_share * time_taper)
    actual_release = float(min(scheduled_spr, actual_release))
    return actual_release, actual_release / scheduled_spr
\end{paperlisting}

\subsection*{F 蒙特卡洛情景预测}

代码~\ref{lst:monte-carlo} 展示长期蒙特卡洛情景树的核心思想。模型以压力指数为主轴，同时扰动供应中断、SPR 释放、绕道能力、需求弹性、风险权重和不确定性溢价，由此得到长期价格区间和尾部风险概率。

\begin{paperlisting}{lst:monte-carlo}{python}{蒙特卡洛情景树参数抽样}
def sample_parameters(rng, base_assumptions, base_behavior, historical_factors):
    raw_stress = float(rng.beta(2.1, 2.3))
    stress = float(np.clip(raw_stress + rng.normal(
        0.0, historical_factors.monte_carlo_stress_noise_sigma), 0.0, 1.0))

    supply_interruption = np.clip(1400 + 400 * stress + rng.normal(0, 35), 1400, 1800)
    spr_max_release = np.clip(700 - 500 * (stress ** 1.08) + rng.normal(0, 30), 200, 700)
    route_max_capacity = np.clip(360 - 210 * stress + rng.normal(0, 22), 150, 420)
    long_elasticity = np.clip(-0.35 + 0.25 * stress + rng.normal(0, 0.018),
                              -0.35, -0.10)

    risk_weight = np.clip(base_behavior.risk_weight * (0.74 + 0.62 * stress
                                                       + rng.normal(0, 0.06)),
                          1.65, 3.35)
    uncertainty_scale = 0.64 + 0.78 * stress + rng.normal(0, 0.06)
    uncertainty_floor = np.clip(base_behavior.uncertainty_floor * uncertainty_scale,
                                0.13, 0.36)

    assumptions = replace(base_assumptions,
                          supply_interruption=supply_interruption,
                          spr_max_release=spr_max_release,
                          route_max_capacity=route_max_capacity,
                          long_elasticity=long_elasticity)
    behavior = replace(base_behavior,
                       risk_weight=risk_weight,
                       uncertainty_floor=uncertainty_floor)
    return assumptions, behavior, {"stress_index": stress}

def run_monte_carlo(n_samples):
    forecast_frame, prefix, base_assumptions, base_behavior = load_baseline()
    historical_factors = load_historical_model_factors()
    rng = np.random.default_rng(RANDOM_SEED)
    paths, metric_rows = [], []

    for sample_id in range(1, n_samples + 1):
        assumptions, behavior, meta = sample_parameters(
            rng, base_assumptions, base_behavior, historical_factors)
        path = simulate_one_sample(forecast_frame, prefix, assumptions, behavior)
        paths.append(path)
        metric_rows.append(summarize_sample(path, prefix, meta))

    metrics = pd.DataFrame(metric_rows)
    return metrics, build_quantiles(paths), build_tail_risk(metrics)
\end{paperlisting}

\subsection*{G 机制消融与统计检验}

代码~\ref{lst:ablation} 用于逐项关闭模型机制，检验主结论是否依赖某一个单独设定。若去除风险溢价、缓冲折价或预期修复后误差显著上升，则说明这些机制对解释价格平台具有实质贡献。

\begin{paperlisting}{lst:ablation}{python}{机制消融实验设计}
VARIANTS = [
    ("完整统一模型", "题面物理机制 + 市场价格形成机制", base_variant),
    ("去除SPR释放", "关闭战略石油储备释放", no_spr),
    ("去除商业库存缓冲", "关闭库存吸收缺口能力", no_inventory),
    ("去除绕道运输", "关闭绕道运输恢复能力", no_reroute),
    ("去除需求收缩", "关闭高油价需求下降", no_demand_contraction),
    ("去除风险与不确定性溢价", "关闭地缘风险和不确定性平台", no_risk_uncertainty),
    ("去除缓冲确认折价", "关闭市场确认缓冲后的降温", no_buffer_confirmation),
    ("去除预期修复折价", "关闭中后期预期修复再定价", no_expectation_relief),
]

def run_ablation(event_df, base_assumptions, base_behavior):
    rows = []
    for name, note, variant_fn in VARIANTS:
        assumptions, behavior = variant_fn(base_assumptions, base_behavior)
        simulation = dynamic.simulate_dynamic_model(event_df, assumptions, behavior)
        metrics = evaluate_simulation(simulation)
        rows.append({
            "消融实验": name,
            "说明": note,
            "RMSE": metrics["RMSE"],
            "MAE": metrics["MAE"],
            "峰值误差": metrics["峰值误差"],
        })
    return pd.DataFrame(rows)
\end{paperlisting}

\subsection*{H 关键结果核查表}

\begin{table}[htbp]
\centering
\caption*{附表 1：关键结果核查}
\begin{tabularx}{\linewidth}{lX}
\toprule
项目 & 核查结果 \\
\midrule
真实数据来源 & 所有真实价格均来自赛题附件 CSV，外部数据只用于合理性校验 \\
短期窗口 & 2026 年 3 月 3 日至 2026 年 5 月 5 日，共 46 个交易日 \\
短期拟合误差 & RMSE = 3.44 美元/桶，MAE = 2.85 美元/桶，MAPE = 2.86\% \\
峰值误差 & 模拟峰值 110.42 美元/桶，实际最高收盘价 114.06 美元/桶 \\
基准对比 & 主模型优于朴素上一日基准和滚动 ARIMA 基准 \\
长期中性结果 & 第 180 天价格约 94.15 美元/桶 \\
长期悲观结果 & 第 180 天价格约 110.16 美元/桶，外推期保留二次跳涨风险 \\
尾部风险 & 蒙特卡洛显示外推期最高价突破 120 美元/桶的条件概率约 16.3\% \\
最敏感因素 & 封锁风险衰减速度、地缘风险权重、不确定性与制度风险强度 \\
\bottomrule
\end{tabularx}
\end{table}

\subsection*{I 一键构建与独立审计补充}

代码~\ref{lst:build-pipeline} 展示最终论文的自动构建入口。该脚本将数据清洗、外部数据校验、情景预测、敏感性分析、蒙特卡洛、统计审计、图表生成和 PDF/DOCX 导出串联为统一流水线，用于避免手工复制结果造成口径不一致。

\begin{paperlisting}{lst:build-pipeline}{bash}{最终论文一键复现脚本核心流程}
#!/usr/bin/env bash
set -euo pipefail

ROOT_DIR="$(cd "$(dirname "${BASH_SOURCE[0]}")/../.." && pwd)"
cd "${ROOT_DIR}"

if [ -x /opt/homebrew/Caskroom/miniconda/base/envs/mathmodel-oil/bin/python3 ]; then
  PYTHON_BIN="/opt/homebrew/Caskroom/miniconda/base/envs/mathmodel-oil/bin/python3"
else
  PYTHON_BIN="python3"
fi

require_artifact "output/calibration/动态模型校准后路径.csv" "Run FULL_REBUILD=1 first."

"${PYTHON_BIN}" -m src.pipeline.clean_data
"${PYTHON_BIN}" -m src.scenarios.scenario_forecast
"${PYTHON_BIN}" -m src.analysis.sensitivity_analysis
"${PYTHON_BIN}" -m src.analysis.monte_carlo_scenario_tree
"${PYTHON_BIN}" -m src.analysis.statistical_audit
"${PYTHON_BIN}" -m src.visualization.final_paper_figures

if command -v Rscript >/dev/null 2>&1; then
  Rscript scripts/audit/r_econometric_audit.R || true
fi

xelatex -shell-escape -interaction=nonstopmode -halt-on-error -output-directory=output/build/latex paper/总论文.tex
pandoc --from=latex --to=docx --toc --output=output/final/论文.docx paper/总论文.tex
\end{paperlisting}

代码~\ref{lst:r-audit} 给出 R 语言独立计量审计的核心逻辑。Python 负责机制递推和情景生成，R 侧用于复核主模型相对朴素基准的误差优势、残差相关性和稳健标准误结果，从而降低单一语言实现自证的风险。

\begin{paperlisting}{lst:r-audit}{r}{R 语言计量审计核心代码}
model_df <- read_csv(calibrated_path, show_col_types = FALSE) %>%
  arrange(trade_date) %>%
  mutate(
    naive_forecast = lag(actual_price),
    model_error = simulated_price - actual_price,
    naive_error = naive_forecast - actual_price
  ) %>%
  filter(!is.na(naive_forecast))

model_error <- model_df$model_error
naive_error <- model_df$naive_error

# Diebold-Mariano 检验：模型损失是否显著低于朴素上一日基准
dm_sq <- forecast::dm.test(model_error, naive_error,
                           alternative = "less", h = 1, power = 2)
dm_abs <- forecast::dm.test(model_error, naive_error,
                            alternative = "less", h = 1, power = 1)

# 校准回归与 Newey-West 稳健标准误
calibration_lm <- lm(actual_price ~ simulated_price, data = model_df)
calibration_hac <- lmtest::coeftest(
  calibration_lm,
  vcov. = sandwich::NeweyWest(calibration_lm, lag = 3,
                              prewhite = FALSE, adjust = TRUE)
)

# 残差自相关、条件异方差和波动聚集检验
lb_resid_5 <- Box.test(model_error, lag = 5, type = "Ljung-Box")
lb_sq_5 <- Box.test(model_error^2, lag = 5, type = "Ljung-Box")
arch_test <- FinTS::ArchTest(model_error, lags = 5)
\end{paperlisting}

代码~\ref{lst:sensitivity-specs} 展示敏感性分析的扰动设置。扰动对象覆盖题面物理层、政策缓冲层、需求层、行为层和长期机制透明层，用于说明本文的稳健性检验不是只改变单个便利参数，而是围绕主要不确定来源系统展开。

\begin{paperlisting}{lst:sensitivity-specs}{python}{敏感性分析扰动参数设置}
SENSITIVITY_SPECS = [
    {"key": "supply_interruption", "参数": "供应中断量",
     "层级": "题面物理层", "扰动": [1400, 1493, 1600, 1800],
     "解释": "封锁实际造成的供应缺口。"},
    {"key": "spr_max_release", "参数": "SPR释放上限",
     "层级": "题面政策缓冲层", "扰动": [200, 450, 670, 700],
     "解释": "战略石油储备释放能力。"},
    {"key": "route_max_capacity", "参数": "绕道运输能力",
     "层级": "题面物理缓冲层", "扰动": [150, 237, 300, 420],
     "解释": "绕道和替代航线形成的额外供给能力。"},
    {"key": "long_elasticity", "参数": "长期需求弹性",
     "层级": "题面需求层", "扰动": [-0.35, -0.25, -0.18, -0.10],
     "解释": "高油价持续后需求收缩强弱。"},
    {"key": "risk_weight", "参数": "地缘风险权重",
     "层级": "市场价格形成层", "扰动": [1.85, 2.22, 2.46, 2.95],
     "解释": "市场定价地缘风险溢价的强度。"},
    {"key": "uncertainty_floor", "参数": "制度风险强度",
     "层级": "长期机制透明层", "扰动": [0.17, 0.22, 0.247, 0.30],
     "解释": "持续封锁和政策协调失败带来的不确定性。"},
]
\end{paperlisting}

\subsection*{J 程序材料说明}

本文不在正文附录中展开完整 Python 或 R 源代码。原因是完整源程序较长，直接粘贴会削弱论文主体的可读性；同时，数学建模论文更应突出模型假设、递推公式、参数解释和结果检验。若比赛提交系统要求源程序，完整代码可作为独立附件提交，正文附录仅保留复现路径、核心算法代码和关键结果核查。
